\documentclass[11pt]{amsart}
\usepackage{amsmath,amssymb,amsthm,booktabs}
\usepackage[margin=1.15in]{geometry}
\usepackage{hyperref}
\newtheorem{theorem}{Theorem}
\newtheorem{proposition}[theorem]{Proposition}
\newtheorem{lemma}[theorem]{Lemma}
\newtheorem{corollary}[theorem]{Corollary}
\theoremstyle{definition}
\newtheorem{remark}[theorem]{Remark}
\newcommand{\Prob}{\mathbb{P}}
\newcommand{\E}{\mathbb{E}}
\newcommand{\R}{\mathbb{R}}
\newcommand{\Mcl}{M_{\mathrm{c}}}
\newcommand{\Mfr}{M_{\mathrm{f}}}
\newcommand{\Var}{\operatorname{Var}}

\title[The mode of classical and free Brownian motion]{The mode of classical and free Brownian motion\\
with an initial distribution at large time}
\author{Yitzchak Shmalo}
\date{November 2026}

\begin{document}
\begin{abstract}
Let $\mu$ be a probability measure on the real line with mean $m$ and finite variance $v$. We show that the density
of $\mu*N(0,t)$ has a unique global maximizer for $t\ge 600v$ and that the density of the free convolution
$\mu\boxplus S(0,t)$ with the semicircle law of variance $t$ has a unique global maximizer for
$t>7(7+4\sqrt2)v$. Both maximizers converge to $m$ as $t\to\infty$, with explicit bounds of order $t^{-1/2}$, and of
order $t^{-1}$ when the third absolute moment is finite. If the seventh moment is finite, the maximizers are
\[
m-\frac{\mu_3}{2t}+\frac{\mu_5-6v\mu_3}{8t^2}+O(t^{-3})\qquad\text{and}\qquad
m-\frac{3\mu_3}{t}+\frac{5\mu_5-22v\mu_3}{t^2}+O(t^{-3})
\]
respectively, where $\mu_k$ is the $k$th central moment of $\mu$. This gives an answer to a problem of Hasebe and
Ueda on the position of the mode of classical and free Brownian motion with an initial distribution.
\end{abstract}
\maketitle

\section{Introduction}

Let $\mu$ be a probability measure on $\R$. The law at time $t$ of a Brownian motion with initial distribution $\mu$
is the convolution $\mu*N(0,t)$, and the law at time $t$ of a free Brownian motion with initial distribution $\mu$
is the free convolution $\mu\boxplus S(0,t)$, where $S(0,t)$ is the semicircle law with mean $0$ and variance $t$,
whose density is $(2\pi t)^{-1}\sqrt{4t-x^2}$ on $[-2\sqrt t,2\sqrt t]$. Free convolution of arbitrary probability
measures on $\R$ is defined by Bercovici and Voiculescu~\cite{BV}, and the measures $\mu\boxplus S(0,t)$ are
described by Biane~\cite{Biane}: each of them has a continuous density.

A probability measure on $\R$ is unimodal if, for some $a\in\R$, it is the sum of an atom at $a$ and a measure with
a density that is non-decreasing on $(-\infty,a)$ and non-increasing on $(a,\infty)$; such a point $a$ is a mode.
Hasebe and Ueda~\cite{HU} study the unimodality of $\mu*N(0,t)$ and $\mu\boxplus S(0,t)$ for large $t$. They prove
that if
\begin{equation}\label{eq:HU}
\alpha:=\int_\R e^{\varepsilon x^2}\,d\mu(x)<\infty\qquad\text{for some }\varepsilon>0,
\end{equation}
then $\mu*N(0,t)$ is unimodal for all $t\ge 36\log(2\alpha)/\varepsilon$ \cite[Theorem~4.3]{HU}, and that if $\mu$
has compact support of diameter $D$, then $\mu\boxplus S(0,t)$ is unimodal for all $t\ge 4D^2$
\cite[Theorem~3.2]{HU}. Neither statement extends to all $\mu$: there are measures $\mu$ for which $\mu*N(0,t)$ is
not unimodal for any $t>0$ \cite[Proposition~4.2]{HU}, and for every Borel function $f\colon\R\to[0,\infty)$ there
is a measure $\mu$ with $\int f\,d\mu<\infty$ for which $\mu\boxplus S(0,t)$ is not unimodal for any $t>0$
\cite[Proposition~3.4]{HU}. In particular, the condition \eqref{eq:HU} does not imply that $\mu\boxplus S(0,t)$ is
eventually unimodal. For unimodal $\mu$ the law $\mu*N(0,t)$ is unimodal for every $t$, because the normal law is
strongly unimodal in the sense of Ibragimov~\cite{Ibr}; the condition \eqref{eq:HU} does not require $\mu$ to be
unimodal. Unimodality in free probability is studied for freely selfdecomposable laws by Hasebe and
Thorbj\o rnsen~\cite{HT}, for free L\'evy processes by Hasebe and Sakuma~\cite{HS}, and, as log-unimodality, for
free positive multiplicative Brownian motion by Hasebe, Ueda and Wang~\cite{HUW}. The paper~\cite{HU} closes its section on the classical case with the following problem.

\medskip
\noindent\textbf{Problem} (\cite[Problem~4.6]{HU})\textbf{.} \emph{Estimate the position of the mode of classical
Brownian motion with initial distributions satisfying the assumption \eqref{eq:HU}. Our proof shows that for
$t\ge 36\log(2\alpha)/\varepsilon$, the mode is located in the interval $[-\sqrt t/2,\sqrt t/2]$. How about free
Brownian motion?}
\medskip

We consider the global maximizers of the two densities. For a unimodal law with a continuous density that has a
unique global maximizer, this maximizer is the unique mode; for a law that is not unimodal it is still defined, and
it is the natural object in the free case because of \cite[Proposition~3.4]{HU}. A finite variance is enough for
the global maximizer to be unique at large times and to converge to the mean of $\mu$.

Throughout, $X$ is a random variable with law $\mu$ and finite variance, and
\[
m=\E X,\qquad Y=X-m,\qquad v=\Var(X),\qquad \mu_k=\E Y^k,\qquad \beta_k=\E|Y|^k .
\]
If $v=0$, then $\mu*N(0,t)=N(m,t)$ and $\mu\boxplus S(0,t)$ is the semicircle law of variance $t$ centred at $m$, and
both densities have the unique maximizer $m$ for every $t>0$. We therefore assume $v>0$.

\begin{theorem}[classical Brownian motion]\label{thm:classical}
Let $0<v<\infty$ and $t\ge 600v$. The density of $\mu*N(0,t)$ has a unique global maximizer $\Mcl(t)$. It satisfies
\[
|\Mcl(t)-m|\le \frac{\sqrt2\,v}{\sqrt t},
\]
and, if $\beta_3<\infty$,
\[
|\Mcl(t)-m|\le\frac{\beta_3}{t}.
\]
\end{theorem}

\begin{theorem}[free Brownian motion]\label{thm:free}
Let $0<v<\infty$, let $a=7+4\sqrt2$ and let $t>7av$. The density of $\mu\boxplus S(0,t)$ has a unique global
maximizer $\Mfr(t)$. It satisfies
\[
|\Mfr(t)-m|\le \frac{4av}{\sqrt t},
\]
and, if $\beta_3<\infty$,
\[
|\Mfr(t)-m|\le\frac{26\bigl((2+\sqrt2)^3v^{3/2}+\beta_3\bigr)}{t}.
\]
\end{theorem}

Since $v^{3/2}\le\beta_3$ and $(2+\sqrt2)^3=20+14\sqrt2$, the last bound is at most $1061\,\beta_3/t$. The numerical
constants in the two theorems, and the thresholds $600v$ and $7av\approx 88.6\,v$, are not optimal.

\begin{theorem}[expansions]\label{thm:expansion}
Let $0<v<\infty$.
\begin{enumerate}
\item[(a)] If $\beta_3<\infty$, then, as $t\to\infty$,
\[
\Mcl(t)=m-\frac{\mu_3}{2t}+o(t^{-1}),\qquad \Mfr(t)=m-\frac{3\mu_3}{t}+o(t^{-1}).
\]
\item[(b)] If $\beta_7<\infty$, then, as $t\to\infty$,
\begin{align}
\Mcl(t)&=m-\frac{\mu_3}{2t}+\frac{\mu_5-6v\mu_3}{8t^2}+O(t^{-3}),\label{eq:exp-c}\\
\Mfr(t)&=m-\frac{3\mu_3}{t}+\frac{5\mu_5-22v\mu_3}{t^2}+O(t^{-3}).\label{eq:exp-f}
\end{align}
\end{enumerate}
\end{theorem}

Thus, when $\beta_3<\infty$, both maximizers approach the mean at rate $t^{-1}$; this rate is exact whenever
$\mu_3\ne0$, and then to leading order the displacement of the free maximizer from the mean is six times that of
the classical one. If $\mu$ is
symmetric about $m$, then both densities are symmetric about $m$, and the uniqueness in
Theorems~\ref{thm:classical} and~\ref{thm:free} gives $\Mcl(t)=m$ for $t\ge600v$ and $\Mfr(t)=m$ for $t>7av$.

The expansions take a simple form in terms of cumulants. Let $\kappa_n$ denote the cumulants of $\mu*N(0,t)$, so that
$\kappa_2=t+v$, $\kappa_3=\mu_3$ and $\kappa_5=\mu_5-10v\mu_3$. Then \eqref{eq:exp-c} reads
\begin{equation}\label{eq:cum-c}
\Mcl(t)=m-\frac{\kappa_3}{2\kappa_2}+\frac{\kappa_5}{8\kappa_2^2}+O(t^{-3}).
\end{equation}
The first term is the classical relation between the mean, the mode and the third cumulant of a nearly normal
distribution; see Haldane~\cite{Haldane} for distributions with given cumulants and Hall~\cite{Hall} for the mode
of a sum of independent random variables. Free cumulants linearize free convolution in the same way as cumulants
linearize convolution \cite[Lectures~11 and~12]{NS}. Let $r_n$ denote the free cumulants of $\mu\boxplus S(0,t)$, so
that $r_2=t+v$, $r_3=\mu_3$ and $r_5=\mu_5-5v\mu_3$. Then \eqref{eq:exp-f} reads
\begin{equation}\label{eq:cum-f}
\Mfr(t)=m-\frac{3r_3}{r_2}+\frac{5r_5}{r_2^2}+O(t^{-3}).
\end{equation}
Asymptotic expansions in the free central limit theorem, for normalized sums of free identically distributed random
variables, are due to Chistyakov and G\"otze~\cite{CG}. Here the setting is a fixed law perturbed by a semicircular
variable of large variance, and \eqref{eq:exp-f} is derived directly from Biane's description of the density.

Under the assumptions of~\cite{HU} the theorems give the following.

\begin{corollary}\label{cor:HU}
Let $\mu$ satisfy \eqref{eq:HU} and not be a point mass, and put
\[
L=\frac{\log\alpha}{\varepsilon},\qquad Q=8\alpha\Bigl(\frac{3}{2e\varepsilon}\Bigr)^{3/2}.
\]
Then $v\le L$, $|m|\le\sqrt L$ and $\beta_3\le Q$. Consequently:
\begin{enumerate}
\item[(i)] for $t\ge\max\{600L,\,36\log(2\alpha)/\varepsilon\}$, and also for $t\ge72\log(2\alpha^2)/\varepsilon$,
the law $\mu*N(0,t)$ is unimodal, its mode is unique, and the mode lies in $[m-Q/t,\,m+Q/t]$;
\item[(ii)] for $t>7(7+4\sqrt2)L$ the density of $\mu\boxplus S(0,t)$ has a unique global maximizer, which lies
within $26\bigl((2+\sqrt2)^3L^{3/2}+Q\bigr)/t$ of $m$;
\item[(iii)] the expansions \eqref{eq:exp-c} and \eqref{eq:exp-f} hold.
\end{enumerate}
\end{corollary}

\begin{corollary}\label{cor:compact}
Let $\mu$ have compact support of diameter $D>0$. For $t\ge 23D^2$ the law $\mu\boxplus S(0,t)$ is unimodal, its mode
is unique and lies within $136\,D^3/t$ of $m$, and the mode satisfies \eqref{eq:exp-f}.
\end{corollary}

Theorem~\ref{thm:classical} is proved in Section~\ref{sec:classical} and Theorem~\ref{thm:free} in
Section~\ref{sec:free}. Section~\ref{sec:expansion} contains the proof of Theorem~\ref{thm:expansion},
Section~\ref{sec:cor} the proofs of the corollaries, and Section~\ref{sec:numerics} an example and numerical checks.

\section{The classical case}\label{sec:classical}

Replacing $X$ by $X-m$ translates both $\mu*N(0,t)$ and $\mu\boxplus S(0,t)$ by $-m$, so in the proofs of
Theorems~\ref{thm:classical}, \ref{thm:free} and~\ref{thm:expansion} we assume $m=0$; then $Y=X$.

The density of $\mu*N(0,t)$ at $h$ is $(2\pi t)^{-1/2}g(h)$ with
\[
g(h)=\E\exp\Bigl(-\frac{(Y-h)^2}{2t}\Bigr).
\]
The function $g$ is positive and continuously differentiable, and $g'(h)=H(h)/t$, where
\[
H(h)=\E\Bigl[(Y-h)\exp\Bigl(-\frac{(Y-h)^2}{2t}\Bigr)\Bigr].
\]
The function $H$ is continuously differentiable, with
\begin{equation}\label{eq:Hprime}
-H'(h)=\E\Bigl[\Bigl(1-\frac{(Y-h)^2}{t}\Bigr)\exp\Bigl(-\frac{(Y-h)^2}{2t}\Bigr)\Bigr];
\end{equation}
the differentiations under the expectation are justified because the integrands and their derivatives in $h$ are
bounded.

\begin{lemma}\label{lem:c-window}
Let $t\ge600v$. Then $g(h)<g(0)$ for every $h$ with $|h|>\sqrt t/2$.
\end{lemma}

\begin{proof}
By Jensen's inequality, $g(0)\ge\exp(-v/(2t))\ge e^{-1/1200}$. Let $|h|>\sqrt t/2$. On the event
$\{|Y-h|\ge\sqrt t/4\}$ the exponential in the definition of $g(h)$ is at most $e^{-1/32}$. On the complementary
event, $|Y|>|h|-\sqrt t/4>\sqrt t/4$, and by Chebyshev's inequality this has probability at most $16v/t\le 16/600$.
Hence
\[
g(h)\le e^{-1/32}+\frac{16}{600}<0.9960<e^{-1/1200}\le g(0).\qedhere
\]
\end{proof}

\begin{lemma}\label{lem:c-slope}
Let $t\ge 600v$. Then $-H'(h)>1/2$ for every $h$ with $|h|\le\sqrt t/2$.
\end{lemma}

\begin{proof}
For $z\ge0$ we have $(1-z)e^{-z/2}\ge1-3z/2$: if $z\le1$, this follows from $e^{-z/2}\ge1-z/2$ and
$(1-z)(1-z/2)\ge1-3z/2$; if $z>1$, then $(1-z)e^{-z/2}\ge1-z\ge1-3z/2$. Applying this to $z=(Y-h)^2/t$
in \eqref{eq:Hprime} and using $\E(Y-h)^2=v+h^2$ gives
\[
-H'(h)\ge1-\frac{3(v+h^2)}{2t}\ge1-\frac{3}{8}-\frac{3v}{2t}\ge\frac58-\frac1{400}>\frac12 .\qedhere
\]
\end{proof}

\begin{proof}[Proof of Theorem~\ref{thm:classical}]
Since $g$ is continuous and $g(h)\to0$ as $|h|\to\infty$ by dominated convergence, $g$ attains its maximum. By
Lemma~\ref{lem:c-window} every global maximizer lies in $[-\sqrt t/2,\sqrt t/2]$, and it is a zero of $H$. By
Lemma~\ref{lem:c-slope}, $H$ is strictly decreasing on this interval, so it has at most one zero there. Hence the
global maximizer is unique; denote it by $h_t$. By Lemma~\ref{lem:c-slope} and the mean value theorem,
\begin{equation}\label{eq:c-mvt}
|h_t|\le2\,|H(0)-H(h_t)|=2\,|H(0)|.
\end{equation}
Since $\E Y=0$ and $0\le1-e^{-z}\le\min\{1,z\}$ for $z\ge0$,
\begin{equation}\label{eq:c-H0}
|H(0)|=\Bigl|\E\Bigl[Y\Bigl(1-e^{-Y^2/(2t)}\Bigr)\Bigr]\Bigr|\le\E\Bigl[|Y|\min\Bigl\{1,\frac{Y^2}{2t}\Bigr\}\Bigr].
\end{equation}
The minimum is at most $|Y|/\sqrt{2t}$, which gives $|H(0)|\le v/\sqrt{2t}$ and $|h_t|\le\sqrt2\,v/\sqrt t$. It is
also at most $Y^2/(2t)$, which gives $|H(0)|\le\beta_3/(2t)$ and $|h_t|\le\beta_3/t$.
\end{proof}

\section{The free case}\label{sec:free}

We recall the description of $\mu\boxplus S(0,t)$ due to Biane~\cite{Biane}, in the form given
in~\cite[Section~2.2]{HU}. For $u\in\R$ and $R>0$ let
\[
F(u,R)=\E\Bigl[\frac{1}{R+(u-X)^2}\Bigr],
\]
and let
\[
v_t(u)=\inf\Bigl\{w\ge0:\ \E\Bigl[\frac1{w^2+(u-X)^2}\Bigr]\le\frac1t\Bigr\}.
\]
The function $v_t$ is continuous on $\R$, and whenever $v_t(u)>0$ the number $R=v_t(u)^2$ is the unique positive
solution of $F(u,R)=1/t$ \cite[(2.8)--(2.10)]{HU}. The map
\[
\psi_t(u)=u+t\,\E\Bigl[\frac{u-X}{(u-X)^2+v_t(u)^2}\Bigr]
\]
is a homeomorphism of $\R$ onto $\R$, the measure $\mu\boxplus S(0,t)$ has a continuous density $p_t$, and
\begin{equation}\label{eq:biane}
p_t(\psi_t(u))=\frac{v_t(u)}{\pi t},\qquad u\in\R
\end{equation}
\cite[(2.13) and (2.14)]{HU}. Consequently the global maximizers of $p_t$ are the images under $\psi_t$ of the global
maximizers of $v_t$.

For each $u$ the function $R\mapsto F(u,R)$ is continuous and strictly decreasing on $(0,\infty)$, and
$F(u,R)\le1/R$. As in Section~\ref{sec:classical} we assume $m=0$, so that $Y=X$.

\begin{lemma}\label{lem:f-max}
Let $t>v$. The function $v_t$ attains its maximum, and $R^*:=\max_u v_t(u)^2$ satisfies $t-v\le R^*\le t$. The set of
global maximizers of $v_t$ is $\{u\in\R:F(u,R^*)=1/t\}$, and $F(u,R^*)\le1/t$ for every $u\in\R$.
\end{lemma}

\begin{proof}
The function $x\mapsto(R+x)^{-1}$ is convex on $[0,\infty)$, so Jensen's inequality gives
$F(0,t-v)\ge(t-v+\E Y^2)^{-1}=1/t$. Hence $v_t(0)>0$, and since $F(0,\cdot)$ is decreasing, $v_t(0)^2\ge t-v$. For
fixed $\delta>0$, dominated convergence gives $F(u,\delta^2)\to0$ as $|u|\to\infty$, so $v_t(u)\le\delta$ for all
large $|u|$. Thus the continuous function $v_t$ tends to $0$ at infinity and is positive at $0$, and it attains its
maximum, with $R^*\ge t-v$. At a maximizer $u^*$ we have $1/t=F(u^*,R^*)\le1/R^*$, so $R^*\le t$.

Let $u\in\R$. If $v_t(u)>0$, then $v_t(u)^2\le R^*$ and $F(u,R^*)\le F(u,v_t(u)^2)=1/t$, with equality if and only if
$v_t(u)^2=R^*$. If $v_t(u)=0$, then $F(u,w^2)\le1/t$ for every $w>0$, and $F(u,R^*)<1/t$ because $F(u,\cdot)$ is
strictly decreasing. This proves the last two assertions.
\end{proof}

Let $B=(2+\sqrt2)\sqrt v$, so that $B^2+v=av$ with $a=7+4\sqrt2$.

\begin{lemma}\label{lem:f-window}
Let $t\ge4v$. Every global maximizer $u^*$ of $v_t$ satisfies $|u^*|\le B$.
\end{lemma}

\begin{proof}
Since $F(u^*,R^*)=1/t$ and $R^*\ge t-v$,
\begin{equation}\label{eq:f-mass}
\E\Bigl[\frac{(u^*-Y)^2}{R^*+(u^*-Y)^2}\Bigr]=1-R^*F(u^*,R^*)=1-\frac{R^*}{t}\le\frac vt .
\end{equation}
By Chebyshev's inequality, $\Prob(|Y|\le\sqrt{2v})\ge1/2$. Suppose that $|u^*|>\sqrt{2v}$ and put
$b=|u^*|-\sqrt{2v}$. On the event $\{|Y|\le\sqrt{2v}\}$ we have $|u^*-Y|\ge b$, and since $x\mapsto x/(R^*+x)$ is
increasing and $R^*\le t$, the integrand in \eqref{eq:f-mass} is at least $b^2/(t+b^2)$ there. Hence
$b^2/(2(t+b^2))\le v/t$, that is, $b^2(t-2v)\le2vt$. For $t\ge4v$ this gives $b^2\le4v$ and
$|u^*|\le(2+\sqrt2)\sqrt v$.
\end{proof}

\begin{lemma}\label{lem:f-concave}
Let $t>7av$. The function $u\mapsto F(u,R^*)$ is strictly concave on $[-B,B]$.
\end{lemma}

\begin{proof}
For fixed $R>0$ the function $u\mapsto(R+(u-x)^2)^{-1}$ and its first two derivatives are bounded uniformly in $x$,
so $F(\cdot,R)$ is twice continuously differentiable and
\[
\partial_u^2F(u,R)=-2\,\E\Bigl[\frac{R-3(u-Y)^2}{(R+(u-Y)^2)^3}\Bigr].
\]
For $z\ge0$,
\[
\frac{1-3z}{(1+z)^3}-(1-6z)=\frac{z^2(15+17z+6z^2)}{(1+z)^3}\ge0 .
\]
With $z=(u-Y)^2/R$ and $\E(u-Y)^2=v+u^2\le av$ for $|u|\le B$ this gives
\[
\E\Bigl[\frac{R-3(u-Y)^2}{(R+(u-Y)^2)^3}\Bigr]\ge\frac1{R^2}\Bigl(1-\frac{6av}{R}\Bigr).
\]
If $t>7av$, then $R^*\ge t-v>6av$ by Lemma~\ref{lem:f-max}, and $\partial_u^2F(u,R^*)<0$ for $|u|\le B$.
\end{proof}

\begin{proof}[Proof of Theorem~\ref{thm:free}]
Let $t>7av$. By Lemma~\ref{lem:f-max} the global maximizers of $v_t$ are the points at which $F(\cdot,R^*)$ attains
its maximum value $1/t$; by Lemma~\ref{lem:f-window} they all lie in $[-B,B]$; and by Lemma~\ref{lem:f-concave} the
function $F(\cdot,R^*)$ has at most one maximizer in this interval. Hence $v_t$ has a unique global maximizer $u^*$,
and $p_t$ has the unique global maximizer $\Mfr(t)=\psi_t(u^*)$.

Write $q=u^*-Y$. Since $u^*$ maximizes the differentiable function $F(\cdot,R^*)$, we have
$\partial_uF(u^*,R^*)=0$, that is,
\begin{equation}\label{eq:f-stat}
\E\Bigl[\frac{q}{(1+q^2/R^*)^2}\Bigr]=0 .
\end{equation}
As $\E q=u^*$, this gives
\begin{equation}\label{eq:f-u}
u^*=\E\Bigl[q\Bigl(1-\frac1{(1+q^2/R^*)^2}\Bigr)\Bigr].
\end{equation}
Moreover, by the definition of $\psi_t$,
\begin{equation}\label{eq:f-M}
\Mfr(t)=u^*+\frac{t}{R^*}\,\E\Bigl[\frac{q}{1+q^2/R^*}\Bigr]
=u^*+\frac{t}{R^*}\Bigl(u^*-\E\Bigl[\frac{q^3}{R^*+q^2}\Bigr]\Bigr).
\end{equation}
Since $v<t/7$, Lemma~\ref{lem:f-max} gives $R^*\ge6t/7$, and $\E q^2=v+(u^*)^2\le av$ by Lemma~\ref{lem:f-window}.

For $z\ge0$ we have $0\le1-(1+z)^{-2}\le\min\{1,2z\}\le\sqrt{2z}$. Hence \eqref{eq:f-u} gives
$|u^*|\le\sqrt2\,\E q^2/\sqrt{R^*}\le\sqrt2\,av/\sqrt{R^*}$. Since $R^*+q^2\ge2\sqrt{R^*}|q|$, we also have
$|\E[q^3/(R^*+q^2)]|\le\E q^2/(2\sqrt{R^*})\le av/(2\sqrt{R^*})$. By \eqref{eq:f-M},
\[
|\Mfr(t)|\le\Bigl(1+\frac76\Bigr)\frac{\sqrt2\,av}{\sqrt{R^*}}+\frac76\cdot\frac{av}{2\sqrt{R^*}}
\le\Bigl(\frac{13\sqrt2}{6}+\frac7{12}\Bigr)\sqrt{\frac76}\,\frac{av}{\sqrt t}<\frac{4av}{\sqrt t}.
\]

Let now $\beta_3<\infty$. From $|x-y|^3\le4(|x|^3+|y|^3)$ and $|u^*|\le B$ we get $\E|q|^3\le K:=4(B^3+\beta_3)$.
Using $1-(1+z)^{-2}\le2z$ in \eqref{eq:f-u} gives
\begin{equation}\label{eq:f-u3}
|u^*|\le\frac{2K}{R^*},
\end{equation}
and $|\E[q^3/(R^*+q^2)]|\le K/R^*$. By \eqref{eq:f-M},
\[
|\Mfr(t)|\le\Bigl(1+\frac76\Bigr)\frac{2K}{R^*}+\frac76\cdot\frac{K}{R^*}=\frac{11K}{2R^*}\le\frac{77K}{12\,t}
<\frac{26(B^3+\beta_3)}{t}.\qedhere
\]
\end{proof}

\begin{remark}
The proof shows that the global maximizer of $p_t$ is unique for $t>7av$ whether or not $\mu\boxplus S(0,t)$ is
unimodal. The density may have further local maxima at such times; an example is given in
Section~\ref{sec:numerics}.
\end{remark}

\section{Expansions}\label{sec:expansion}

We keep the assumption $m=0$, assume $\beta_3<\infty$, as in both parts of Theorem~\ref{thm:expansion}, and write
$s=1/t$. In this section $O(\cdot)$ and $o(\cdot)$ refer to $s\to0$, with
constants depending on $\mu$.

\subsection{The classical maximizer}
Let $h=h_t$ be the maximizer of Section~\ref{sec:classical} and put $q=Y-h$. By Theorem~\ref{thm:classical},
$|h|\le\beta_3s$, so $|q|\le|Y|+1$ for small $s$. Since $\E Y=0$,
\begin{equation}\label{eq:c-moments}
\E q=-h,\qquad \E q^3=\mu_3-3vh-h^3,\qquad \E q^5=\mu_5-5\mu_4h+10\mu_3h^2-10vh^3-h^5,
\end{equation}
where the last identity requires $\beta_5<\infty$. The equation $H(h)=0$ reads $\E[q\,e^{-sq^2/2}]=0$.

\emph{Proof of {\rm(a)}.} For $x\ge0$ write $e^{-x}=1-x+\rho(x)$ with $0\le\rho(x)\le\min\{x,x^2/2\}$. Then
\[
0=-h-\frac s2\,\E q^3+\E\bigl[q\,\rho(sq^2/2)\bigr],\qquad
\bigl|\E\bigl[q\,\rho(sq^2/2)\bigr]\bigr|\le s\,\E\Bigl[|q|^3\min\Bigl\{\frac12,\frac{sq^2}{8}\Bigr\}\Bigr].
\]
The last integrand is at most $(|Y|+1)^3/2$, which is integrable, and it tends to $0$ pointwise as $s\to0$; by
dominated convergence the last expectation is $o(1)$. Since $\E q^3=\mu_3+O(s)$, we get $h=-\mu_3s/2+o(s)$.

\emph{Proof of {\rm(b)}.} Let $\beta_7<\infty$. For $x\ge0$, $e^{-x}=1-x+x^2/2+\rho_3(x)$ with $|\rho_3(x)|\le x^3/6$.
Since $\E|q|^7$ is bounded for small $s$, the equation $\E[q\,e^{-sq^2/2}]=0$ gives
\[
0=\E q-\frac s2\,\E q^3+\frac{s^2}8\,\E q^5+O(s^3).
\]
Inserting \eqref{eq:c-moments} and $h=O(s)$,
\[
h=-\frac s2\,(\mu_3-3vh)+\frac{s^2}{8}\,\mu_5+O(s^3).
\]
Hence $h=-\mu_3s/2+O(s^2)$, and substituting this into the right-hand side,
\[
h=-\frac{\mu_3}{2}\,s+\frac{\mu_5-6v\mu_3}{8}\,s^2+O(s^3),
\]
which is \eqref{eq:exp-c}.

\subsection{The free maximizer}
Let $t>7av$, let $u=u^*$ and $R^*$ be as in Section~\ref{sec:free}, and put $r=R^*/t$ and $q=u-Y$. By
Lemma~\ref{lem:f-max}, $1-vs\le r\le1$, and by \eqref{eq:f-u3}, $u=O(s)$. Since $|q|\le|Y|+B$, the moments
$\E|q|^k$ with $k\le7$ are bounded whenever $\beta_k<\infty$. Since $\E Y=0$,
\begin{equation}\label{eq:f-moments}
\E q=u,\quad \E q^2=v+u^2,\quad \E q^3=u^3+3vu-\mu_3,\quad \E q^5=u^5+10vu^3-10\mu_3u^2+5\mu_4u-\mu_5,
\end{equation}
where the last identity requires $\beta_5<\infty$, and $\E q^4=\mu_4+O(s)$ if $\beta_4<\infty$. In terms of $r$ and
$s$, the equation $F(u,R^*)=1/t$, the stationarity
condition \eqref{eq:f-stat} and the formula \eqref{eq:f-M} read
\begin{equation}\label{eq:f-system}
\E\Bigl[\frac1{r+sq^2}\Bigr]=1,\qquad \E\Bigl[\frac{q}{(r+sq^2)^2}\Bigr]=0,\qquad
\Mfr(t)=u+\E\Bigl[\frac{q}{r+sq^2}\Bigr].
\end{equation}
We use the following expansions, valid for $z\ge0$:
\begin{align}
\frac1{1+z}&=1-z+\rho_1(z)=1-z+z^2+\sigma_1(z), & 0\le\rho_1(z)&\le\min\{z,z^2\}, & |\sigma_1(z)|&\le z^3,
\label{eq:exp1}\\
\frac1{(1+z)^2}&=1-2z+\rho_2(z)=1-2z+3z^2+\sigma_2(z), & 0\le\rho_2(z)&\le\min\{2z,3z^2\}, & |\sigma_2(z)|&\le4z^3.
\label{eq:exp2}
\end{align}
Here $\rho_1(z)=z^2/(1+z)$ and $\sigma_1(z)=-z^3/(1+z)$; the bounds on $\rho_2$ and $\sigma_2$ follow from Taylor's
formula, since the second and third derivatives of $(1+z)^{-2}$ are $6(1+z)^{-4}$ and $-24(1+z)^{-5}$, together with
$(1+z)^{-2}\le1$. Throughout we take $z=sq^2/r$.

\emph{Proof of {\rm(a)}.} Multiplying the second equation in \eqref{eq:f-system} by $r^2$ and
using \eqref{eq:exp2},
\[
0=u-\frac{2s}{r}\,\E q^3+\E\bigl[q\,\rho_2(z)\bigr],\qquad
\bigl|\E\bigl[q\,\rho_2(z)\bigr]\bigr|\le\frac sr\,\E\Bigl[|q|^3\min\Bigl\{2,\frac{3sq^2}{r}\Bigr\}\Bigr]=o(s),
\]
by dominated convergence as in the classical case. Since $\E q^3=-\mu_3+O(s)$ and $1/r=1+O(s)$, this gives
$u=-2\mu_3s+o(s)$. In the same way, the third relation in \eqref{eq:f-system} and \eqref{eq:exp1} give
\[
\Mfr(t)=u+\frac1r\Bigl(u-\frac sr\,\E q^3+\E\bigl[q\,\rho_1(z)\bigr]\Bigr)=2u+\mu_3s+o(s)=-3\mu_3s+o(s).
\]

\emph{Proof of {\rm(b)}.} Let $\beta_7<\infty$. Multiplying the first equation in \eqref{eq:f-system} by $r$ and
using \eqref{eq:exp1},
\[
r=1-\frac sr\,\E q^2+\frac{s^2}{r^2}\,\E q^4+O(s^3).
\]
With \eqref{eq:f-moments} this gives $r=1-vs+O(s^2)$ and then
\[
r=1-vs+(\mu_4-v^2)s^2+O(s^3),\qquad \frac1r=1+vs+O(s^2),\qquad \frac1{r^2}=1+2vs+O(s^2).
\]
Multiplying the second equation in \eqref{eq:f-system} by $r^2$ and using \eqref{eq:exp2},
\[
0=\E q-\frac{2s}{r}\,\E q^3+\frac{3s^2}{r^2}\,\E q^5+O(s^3)
=u-6vsu+2\mu_3s+2v\mu_3s^2-3\mu_5s^2+O(s^3),
\]
where we inserted \eqref{eq:f-moments} and $u=O(s)$. Hence $u=-2\mu_3s+O(s^2)$ and then
\[
u=-2\mu_3s+(3\mu_5-14v\mu_3)s^2+O(s^3).
\]
Finally, by the third relation in \eqref{eq:f-system} and \eqref{eq:exp1},
\begin{align*}
\Mfr(t)&=u+\frac1r\,\E q-\frac{s}{r^2}\,\E q^3+\frac{s^2}{r^3}\,\E q^5+O(s^3)\\
&=2u-2vsu+\mu_3s+(2v\mu_3-\mu_5)s^2+O(s^3)\\
&=-3\mu_3s+(5\mu_5-22v\mu_3)s^2+O(s^3),
\end{align*}
which is \eqref{eq:exp-f}. This completes the proof of Theorem~\ref{thm:expansion}.

The forms \eqref{eq:cum-c} and \eqref{eq:cum-f} follow by expanding $1/(t+v)$ and $1/(t+v)^2$ in powers of $s$: the
coefficients of $s^2$ are $v\mu_3/2+(\mu_5-10v\mu_3)/8=(\mu_5-6v\mu_3)/8$ and
$3v\mu_3+5(\mu_5-5v\mu_3)=5\mu_5-22v\mu_3$. The cumulants and free cumulants used there are those of a centred
law: $\kappa_5=\mu_5-10\mu_2\mu_3$, and $r_5=\mu_5-5\mu_2\mu_3$ \cite[Lecture~11]{NS}; the normal law has no
cumulants of order higher than two, and the semicircle law has no free cumulants of order higher than two.

\section{Proofs of the corollaries}\label{sec:cor}

\begin{proof}[Proof of Corollary~\ref{cor:HU}]
By Jensen's inequality, $\exp(\varepsilon\,\E X^2)\le\alpha$, so $v\le\E X^2\le L$ and $m^2\le L$. For $k>0$ the
function $x\mapsto x^ke^{-\varepsilon x^2}$ on $[0,\infty)$ has the maximum $(k/(2e\varepsilon))^{k/2}$, hence
\[
\E|X|^k\le\alpha\Bigl(\frac{k}{2e\varepsilon}\Bigr)^{k/2}.
\]
In particular all moments of $\mu$ are finite, and $\beta_3\le4(\E|X|^3+|m|^3)\le8\,\E|X|^3\le Q$. Since $\mu$ is
not a point mass, $v>0$.

(i) For $t\ge600L$, Theorem~\ref{thm:classical} applies and gives a unique global maximizer $\Mcl(t)$ of the density
with $|\Mcl(t)-m|\le Q/t$. For $t\ge36\log(2\alpha)/\varepsilon$ the law $\mu*N(0,t)$ is unimodal
by \cite[Theorem~4.3]{HU}. The density of a unimodal law is maximal at each of its modes, so the mode is unique and
equal to $\Mcl(t)$.

Now let $t\ge72\log(2\alpha^2)/\varepsilon$, and let $Y=X-m$. Since $Y^2\le2X^2+2m^2$ and
$e^{\varepsilon m^2}\le\alpha$, we have $\E e^{\varepsilon Y^2/2}\le\alpha^2$, so the law of $Y$ satisfies
\eqref{eq:HU} with $\varepsilon/2$ in place of $\varepsilon$ and a constant at most $\alpha^2$ in place of $\alpha$.
By \cite[Theorem~4.3]{HU} and the statement on the mode in \cite[Problem~4.6]{HU}, applied to the law of $Y$, the law
$\mu*N(0,t)$ is unimodal and each of its modes has the form $m+h$ with $|h|\le\sqrt t/2$. The density is maximal at
each mode, so $H(h)=0$. Since $t\ge144L\ge36v$, the first inequality in the proof of Lemma~\ref{lem:c-slope} gives
\[
-H'(h)\ge1-\frac38-\frac{3v}{2t}\ge\frac7{12}\qquad(|h|\le\sqrt t/2),
\]
so $H$ has only one zero in $[-\sqrt t/2,\sqrt t/2]$ and the mode is unique. By the mean value theorem
and \eqref{eq:c-H0}, $|h|\le\frac{12}7|H(0)|\le\frac{6}{7}\beta_3/t\le Q/t$.

(ii) and (iii) follow from Theorems~\ref{thm:free} and~\ref{thm:expansion}, since $7av\le7aL$ and
$v^{3/2}\le L^{3/2}$.
\end{proof}

\begin{proof}[Proof of Corollary~\ref{cor:compact}]
Let $[c-D/2,c+D/2]$ be the smallest interval containing the support of $\mu$. Then $v\le\E(X-c)^2\le D^2/4$, and
$|Y|\le D$ gives $\beta_3\le Dv\le D^3/4$. Hence $7av\le7aD^2/4<23D^2$, and for $t\ge23D^2$ Theorem~\ref{thm:free}
gives a unique global maximizer $\Mfr(t)$ of $p_t$ with
\[
|\Mfr(t)-m|\le\frac{26}{t}\Bigl(\frac{(2+\sqrt2)^3}{8}+\frac14\Bigr)D^3<\frac{136\,D^3}{t}.
\]
For $t\ge4D^2$ the law $\mu\boxplus S(0,t)$ is unimodal by \cite[Theorem~3.2]{HU}, and as in the proof of
Corollary~\ref{cor:HU} its mode is unique and equal to $\Mfr(t)$. All moments of $\mu$ are finite, so
Theorem~\ref{thm:expansion} applies.
\end{proof}

\section{An example and numerical checks}\label{sec:numerics}

Let $\mu=\frac23\delta_{-1}+\frac13\delta_{2}$. Then $m=0$, $v=2$, $\mu_3=2$ and $\mu_5=10$, and
Theorem~\ref{thm:expansion} gives
\[
\Mcl(t)=-\frac1t-\frac{7}{4t^2}+O(t^{-3}),\qquad \Mfr(t)=-\frac6t-\frac{38}{t^2}+O(t^{-3}).
\]
Table~\ref{tab:example} compares these with the maximizers computed from the equations $H(h)=0$
and \eqref{eq:f-system}.

\begin{table}[ht]
\centering
\begin{tabular}{rllll}
\toprule
$t$ & $t\,\Mcl(t)$ & $-1-7/(4t)$ & $t\,\Mfr(t)$ & $-6-38/t$\\
\midrule
$10^4$ & $-1.00017502375$ & $-1.000175$ & $-6.00380137932$ & $-6.0038$\\
$10^5$ & $-1.00001750024$ & $-1.0000175$ & $-6.00038001380$ & $-6.00038$\\
$10^6$ & $-1.00000175000$ & $-1.00000175$ & $-6.00003800014$ & $-6.000038$\\
\bottomrule
\end{tabular}
\caption{The global maximizers for $\mu=\frac23\delta_{-1}+\frac13\delta_2$ and their two-term
expansions.}\label{tab:example}
\end{table}

The program \texttt{code/mode\_checks.py} in the repository of the paper derives the coefficients
in \eqref{eq:exp-c}--\eqref{eq:cum-f} symbolically from the equations $H(h)=0$ and \eqref{eq:f-system}, using
SymPy~\cite{sympy}. With mpmath~\cite{mpmath} it solves the equations $H(h)=0$ and \eqref{eq:f-system} near the mean for five laws
with finitely many atoms, for $t$ between $10^3$ and $10^6$, and for a centred exponential law, for $t$ between
$10^3$ and $10^5$; in each case the difference of the solution from the
two-term expansion, multiplied by $t^3$, approaches a constant. For the five atomic laws it also evaluates the two
densities on a grid at $t=600v$ and at a time just above $7av$, and confirms that the global maximizer is unique and
satisfies the bounds of Theorems~\ref{thm:classical} and~\ref{thm:free}, together with the inequalities of
Lemmas~\ref{lem:c-slope}, \ref{lem:f-max}, \ref{lem:f-window} and~\ref{lem:f-concave}. For the law
$0.8\,\delta_0+0.199\,\delta_1+0.001\,\delta_{50}$, at a time just above $7av$, the free density has a second local
maximum besides its global maximizer, so $\mu\boxplus S(0,t)$ is not unimodal at that time. A second program,
\texttt{code/checks.py}, locates both global maximizers of four laws with finitely many atoms by a grid search
followed by Newton's method at forty-digit precision and compares them with the expansions. The polynomial identity
in the proof of Lemma~\ref{lem:f-concave} is also checked in Lean~4~\cite{Lean4} with Mathlib~\cite{mathlib}
(\texttt{lean/Algebra.lean}).

\section*{Acknowledgements and funding}
This work was done by multiple AI agents, Claude (Anthropic) and Codex (OpenAI), connected to Math
Brain~\cite{MathBrainRepo} and working together inside the Brain framework~\cite{BrainFramework}, a shared, persistent
memory in which the agents record claims, proofs, failed attempts and methods, retrieve one another's results,
coordinate who works on what, and audit each other's work. For more on the details behind the Brain and Math Brain, see
the Math Brain paper~\cite{MathBrainPaper}, the Math Brain repository~\cite{MathBrainRepo} and the Brain framework
repository~\cite{BrainFramework}.

The research presented in this paper was supported by the European Research Council (ERC) under the European Union's
Horizon Europe research and innovation programme (grant agreement No.~101041711), by the Simons Foundation as part of
the Collaboration on the Mathematical and Scientific Foundations of Deep Learning, by Heights Labs, by Convex Nexus
Capital, by the Israel Science Foundation (grant number 2258/19), and by the Israel Science Foundation (ISF Grant
4101/25). Computational resources were provided in part by Google Cloud through its education credits program (credit
530550103).

\section*{Declaration of competing interest}
The author holds equity in Heights Labs and in Convex Nexus Capital and has a management role at Convex Nexus Capital.
These entities are included in the funding declaration above.

\section*{Declaration of generative AI and AI-assisted technologies}
The paper comes out of Math Brain, an ongoing project on the use of AI systems for research mathematics, and its
results were developed with the Brain~\cite{BrainFramework} and Math Brain~\cite{MathBrainRepo}. In October 2026 a
Codex agent (OpenAI) connected to the Brain, working from the project's list of open problems in probability,
obtained the expansions of Theorem~\ref{thm:expansion}{\rm(b)}, the uniqueness of the free maximizer and bounds of
order $t^{-1}$ for both maximizers, together with numerical checks and the Lean file; further AI agents, among them
Claude Opus~5.5 (Anthropic), checked that work and wrote the program \texttt{code/checks.py}. In November 2026 a Claude Fable
agent (Anthropic) obtained Theorems~\ref{thm:classical}, \ref{thm:free} and~\ref{thm:expansion} and the two
corollaries in their present form with their proofs, wrote and ran the programs and prepared the manuscript, and a further Codex agent, working from the manuscript
alone with its own programs, checked the proofs and recomputed the coefficients and the example. These
checks are not peer review. The author directed the work and is responsible for the content of the paper and for
decisions concerning its dissemination.

\section*{Code and data availability}
Programs and outputs:\\
\url{https://github.com/yspennstate/brownian-mode-location}.

The paper is part of the Math Brain project, whose public repository is
\url{https://github.com/yspennstate/math-brain-public}~\cite{MathBrainRepo}.

\begin{thebibliography}{99}
\bibitem{BV} H.~Bercovici and D.~Voiculescu, Free convolution of measures with unbounded support, \emph{Indiana Univ.
Math. J.} \textbf{42} (1993), no.~3, 733--773.
\href{https://doi.org/10.1512/iumj.1993.42.42033}{doi:10.1512/iumj.1993.42.42033}.
\bibitem{Biane} P.~Biane, On the free convolution with a semi-circular distribution, \emph{Indiana Univ. Math. J.}
\textbf{46} (1997), no.~3, 705--718.
\href{https://doi.org/10.1512/iumj.1997.46.1467}{doi:10.1512/iumj.1997.46.1467}.
\bibitem{CG} G.~P.~Chistyakov and F.~G\"otze, Asymptotic expansions in the CLT in free probability, \emph{Probab.
Theory Related Fields} \textbf{157} (2013), 107--156.
\href{https://doi.org/10.1007/s00440-012-0451-2}{doi:10.1007/s00440-012-0451-2}; arXiv:1109.4844.
\bibitem{Haldane} J.~B.~S.~Haldane, The mode and median of a nearly normal distribution with given cumulants,
\emph{Biometrika} \textbf{32} (1942), 294--299.
\href{https://doi.org/10.1093/biomet/32.3-4.294}{doi:10.1093/biomet/32.3-4.294}.
\bibitem{Hall} P.~Hall, On the limiting behaviour of the mode and median of a sum of independent random variables,
\emph{Ann. Probab.} \textbf{8} (1980), 419--430.
\href{https://doi.org/10.1214/aop/1176994717}{doi:10.1214/aop/1176994717}.
\bibitem{HS} T.~Hasebe and N.~Sakuma, Unimodality for free L\'evy processes, \emph{Ann. Inst. Henri Poincar\'e
Probab. Stat.} \textbf{53} (2017), no.~2, 916--936.
\href{https://doi.org/10.1214/16-AIHP742}{doi:10.1214/16-AIHP742}.
\bibitem{HT} T.~Hasebe and S.~Thorbj\o rnsen, Unimodality of the freely selfdecomposable probability laws, \emph{J.
Theoret. Probab.} \textbf{29} (2016), no.~3, 922--940.
\href{https://doi.org/10.1007/s10959-015-0595-y}{doi:10.1007/s10959-015-0595-y}.
\bibitem{HU} T.~Hasebe and Y.~Ueda, Large time unimodality for classical and free Brownian motions with initial
distributions, \emph{ALEA Lat. Am. J. Probab. Math. Stat.} \textbf{15} (2018), 353--374.
\href{https://doi.org/10.30757/ALEA.v15-15}{doi:10.30757/ALEA.v15-15}; arXiv:1710.08240.
\bibitem{HUW} T.~Hasebe, Y.~Ueda and J.-C.~Wang, Log-unimodality for free positive multiplicative Brownian motion,
\emph{Colloq. Math.} \textbf{169} (2022), no.~2, 209--226.
\href{https://doi.org/10.4064/cm8413-6-2021}{doi:10.4064/cm8413-6-2021}.
\bibitem{Ibr} I.~A.~Ibragimov, On the composition of unimodal distributions, \emph{Theory Probab. Appl.} \textbf{1}
(1956), no.~2, 255--260. \href{https://doi.org/10.1137/1101021}{doi:10.1137/1101021}.
\bibitem{mathlib} The mathlib Community, The Lean mathematical library, in \emph{Proceedings of the 9th ACM SIGPLAN
International Conference on Certified Programs and Proofs (CPP 2020)}, ACM, 2020, 367--381.
\href{https://doi.org/10.1145/3372885.3373824}{doi:10.1145/3372885.3373824}.
\bibitem{sympy} A.~Meurer, C.~P.~Smith, M.~Paprocki, et al., SymPy: symbolic computing in Python, \emph{PeerJ Comput.
Sci.} \textbf{3} (2017), e103. \href{https://doi.org/10.7717/peerj-cs.103}{doi:10.7717/peerj-cs.103}.
\bibitem{Lean4} L.~de~Moura and S.~Ullrich, The Lean 4 theorem prover and programming language, in \emph{Automated
Deduction -- CADE 28}, Lecture Notes in Computer Science, Springer, Cham, 2021, 625--635.
\href{https://doi.org/10.1007/978-3-030-79876-5_37}{doi:10.1007/978-3-030-79876-5\_37}.
\bibitem{mpmath} The mpmath development team, \emph{mpmath: a Python library for arbitrary-precision floating-point
arithmetic}, version 1.3.0, 2023, \url{https://mpmath.org/}.
\bibitem{NS} A.~Nica and R.~Speicher, \emph{Lectures on the Combinatorics of Free Probability}, London Mathematical
Society Lecture Note Series \textbf{335}, Cambridge University Press, Cambridge, 2006.
\href{https://doi.org/10.1017/CBO9780511735127}{doi:10.1017/CBO9780511735127}.
\bibitem{MathBrainRepo} Y.~Shmalo, \emph{Math Brain: public snapshot, code and tests}, 2026,
\url{https://github.com/yspennstate/math-brain-public}.
\bibitem{MathBrainPaper} Y.~Shmalo, \emph{Math Brain: shared, auditable memory for multi-agent mathematical research},
October 2026, \url{https://github.com/yspennstate/math-brain-public/tree/main/papers/math-brain-2026-10-24}.
\bibitem{BrainFramework} Y.~Shmalo, \emph{Brain Framework}, public software repository, 2026,
\url{https://github.com/yspennstate/brain-framework-public}.
\end{thebibliography}
\end{document}
